\scanpage{49}
\uline{\unsure{Rmq|Rem}} On a généralisé par la bande la définition de morphisme de générisation entre MD-graphes aux maquettes (pas néc. $0$-connexes ni stables), pour avoir un plus joli énoncé. Les maquettes déduites de $\underline{G}'$ par le procédé décrit s'appellent les \uline{maquettes générisations} de $\underline{G}'$, ou \uline{maquettes quotients}. Elles forment un ens. ordonné, qui est isomorphe, par $\underline{G}\longleftrightarrow C'$, à l'opposé de l'ens ordonné des parties $C'$ de $A'$ (ens des arêtes de $\underline{G}'$), donc, par $\underline{G}\longleftrightarrow B'=A'\setminus C'$, à l'ens ordonné des parties $B'$ de $A'$. Donc si $\mu'=\operatorname{card}A'$, le nb des générisations de $\underline{G}'$ est $2^{\mu'}$.

Appelons \uline{``écart''} d'un morphisme de générisation $\underline{G}'\mathbin{\text{---}}\underline{G}$ de maquettes l'entier naturel $\delta=\operatorname{card}C'$. Alors l'écart \deleted{\unsure{a priori}} \deleted{est compris} \inserted{(pour $\underline{G}'$ fixé)} \deleted{est} peut prendre ttes les valeurs, entre $0$ et $\unsure{\delta_i}(\underline{G}')=\text{codim modulaire de }\underline{G}'$. Il est nul ssi le morphisme de gén. est un \uline{iso}.
\margin{48}
